\section{AMC 12 12B 2004}
        \begin{enumerate}
        	\item (\textit{AMC 12 12B 2004}) At each basketball practice last week, Jenny made twice as many free throws as she made at the previous practice. At her fifth practice she made 48 free throws. How many free throws did she make at the first practice? 


 $(\mathrm {A}) 3\qquad (\mathrm {B}) 6 \qquad (\mathrm {C}) 9 \qquad (\mathrm {D}) 12 \qquad (\mathrm {E}) 15$


	\item (\textit{AMC 12 12B 2004}) In the expression $c\cdot a^b-d$, the values of $a$, $b$, $c$, and $d$ are 0, 1, 2, and 3, although not necessarily in that order. What is the maximum possible value of the result?


 $(\mathrm {A}) 5\qquad (\mathrm {B}) 6 \qquad (\mathrm {C}) 8 \qquad (\mathrm {D}) 9 \qquad (\mathrm {E}) 10$


	\item (\textit{AMC 12 12B 2004}) If $x$ and $y$ are positive integers for which $2^x3^y=1296$, what is the value of $x+y$?


 $(\mathrm {A}) 8\qquad (\mathrm {B}) 9 \qquad (\mathrm {C}) 10 \qquad (\mathrm {D}) 11 \qquad (\mathrm {E}) 12$


	\item (\textit{AMC 12 12B 2004}) An integer $x$, with $10\leq x\leq 99$, is to be chosen. If all choices are equally likely, what is the probability that at least one digit of $x$ is a 7?


 $(\mathrm {A}) \dfrac{1}{9} \qquad (\mathrm {B}) \dfrac{1}{5} \qquad (\mathrm {C}) \dfrac{19}{90} \qquad (\mathrm {D}) \dfrac{2}{9} \qquad (\mathrm {E}) \dfrac{1}{3}$


	\item (\textit{AMC 12 12B 2004}) On a trip from the United States to Canada, Isabella took $d$ U.S. dollars. At the border she exchanged them all, receiving 10 Canadian dollars for every 7 U.S. dollars. After spending 60 Canadian dollars, she had $d$ Canadian dollars left. What is the sum of the digits of $d$? 


 $(\mathrm {A}) 5\qquad (\mathrm {B}) 6 \qquad (\mathrm {C}) 7 \qquad (\mathrm {D}) 8 \qquad (\mathrm {E}) 9$


	\item (\textit{AMC 12 12B 2004}) Minneapolis-St. Paul International Airport is 8 miles southwest of downtown St. Paul and 10 miles southeast of downtown Minneapolis. Which of the following is closest to the number of miles between downtown St. Paul and downtown Minneapolis? 


 $(\mathrm {A}) 13\qquad (\mathrm {B}) 14 \qquad (\mathrm {C}) 15 \qquad (\mathrm {D}) 16 \qquad (\mathrm {E}) 17$


	\item (\textit{AMC 12 12B 2004}) A square has sides of length 10, and a circle centered at one of its vertices has radius 10. What is the area of the union of the regions enclosed by the square and the circle?


 $(\mathrm {A}) 200+25\pi \quad (\mathrm {B}) 100+75\pi \quad (\mathrm {C}) 75+100\pi \quad (\mathrm {D}) 100+100\pi \quad (\mathrm {E}) 100+125\pi$


	\item (\textit{AMC 12 12B 2004}) A grocer makes a display of cans in which the top row has one can and each lower row has two more cans than the row above it. If the display contains 100 cans, how many rows does it contain?


 $(\mathrm {A}) 5 \qquad (\mathrm {B}) 8 \qquad (\mathrm {C}) 9 \qquad (\mathrm {D}) 10 \qquad (\mathrm {E}) 11$


	\item (\textit{AMC 12 12B 2004}) The point $(-3,2)$ is rotated $90^\circ$ clockwise around the origin to point $B$. Point $B$ is then reflected over the line $x=y$ to point $C$. What are the coordinates of $C$?


 $\mathrm{(A)}\ (-3,-2) \qquad \mathrm{(B)}\ (-2,-3) \qquad \mathrm{(C)}\ (2,-3) \qquad \mathrm{(D)}\ (2,3) \qquad \mathrm{(E)}\ (3,2)$


	\item (\textit{AMC 12 12B 2004}) An annulus is the region between two concentric circles. The concentric circles in the figure have radii $b$ and $c$, with $b>c$. Let $OX$ be a radius of the larger circle, let $XZ$ be tangent to the smaller circle at $Z$, and let $OY$ be the radius of the larger circle that contains $Z$. Let $a=XZ$, $d=YZ$, and $e=XY$. What is the area of the annulus? 


 
                    \begin{figure}[H]
                        \centering
                        \includegraphics[width=0.5\textwidth]{Problem/AMC_12_12B_2004_10_Problem_Figure_1.png}
                    \end{figure}
                    
$\mathrm{(A) \ } \pi a^2 \qquad \mathrm{(B) \ } \pi b^2 \qquad \mathrm{(C) \ } \pi c^2 \qquad \mathrm{(D) \ } \pi d^2 \qquad \mathrm{(E) \ } \pi e^2$


	\item (\textit{AMC 12 12B 2004}) All the students in an algebra class took a $100$-point test. Five students scored $100$, each student scored at least $60$, and the mean score was $76$. What is the smallest possible number of students in the class?


 $\mathrm{(A)}\ 10 \qquad \mathrm{(B)}\ 11 \qquad \mathrm{(C)}\ 12 \qquad \mathrm{(D)}\ 13 \qquad \mathrm{(E)}\ 14$


	\item (\textit{AMC 12 12B 2004}) In the sequence $2001$, $2002$, $2003$, $\ldots$ , each term after the third is found by subtracting the previous term from the sum of the two terms that precede that term. For example, the fourth term is $2001 + 2002 - 2003 = 2000$. What is the
$2004^\textrm{th}$ term in this sequence?


 $\mathrm{(A) \ } -2004 \qquad \mathrm{(B) \ } -2 \qquad \mathrm{(C) \ } 0 \qquad \mathrm{(D) \ } 4003 \qquad \mathrm{(E) \ } 6007$


	\item (\textit{AMC 12 12B 2004}) If $f(x) = ax+b$ and $f^{-1}(x) = bx+a$ with $a$ and $b$ real, what is the value of $a+b$? 


 $\mathrm{(A)}\ -2 \qquad\mathrm{(B)}\ -1 \qquad\mathrm{(C)}\ 0 \qquad\mathrm{(D)}\ 1 \qquad\mathrm{(E)}\ 2$


	\item (\textit{AMC 12 12B 2004}) In $\triangle ABC$, $AB=13$, $AC=5$, and $BC=12$. Points $M$ and $N$ lie on $AC$ and $BC$, respectively, with $CM=CN=4$. Points $J$ and $K$ are on $AB$ so that $MJ$ and $NK$ are perpendicular to $AB$. What is the area of pentagon $CMJKN$? 


 
                    \begin{figure}[H]
                        \centering
                        \includegraphics[width=0.5\textwidth]{Problem/AMC_12_12B_2004_14_Problem_Figure_1.png}
                    \end{figure}
                    
$\mathrm{(A)}\ 15 \qquad \mathrm{(B)}\ \frac{81}{5} \qquad \mathrm{(C)}\ \frac{205}{12} \qquad \mathrm{(D)}\ \frac{240}{13} \qquad \mathrm{(E)}\ 20$


	\item (\textit{AMC 12 12B 2004}) The two digits in Jack's age are the same as the digits in Bill's age, but in reverse order. In five years Jack will be twice as old as Bill will be then. What is the difference in their current ages?


 $\mathrm{(A) \ } 9 \qquad \mathrm{(B) \ } 18 \qquad \mathrm{(C) \ } 27 \qquad \mathrm{(D) \ } 36\qquad \mathrm{(E) \ } 45$


	\item (\textit{AMC 12 12B 2004}) A \href{https://artofproblemsolving.com/wiki/index.php?title=Function}{function} $f$ is defined by $f(z) = i\overline{z}$, where $i=\sqrt{-1}$ and $\overline{z}$ is the \href{https://artofproblemsolving.com/wiki/index.php?title=Complex\_conjugate}{complex conjugate} of $z$. How many values of $z$ satisfy both $|z| = 5$ and $f(z) = z$?


 $\mathrm{(A)}\ 0 \qquad\mathrm{(B)}\ 1 \qquad\mathrm{(C)}\ 2  \qquad\mathrm{(D)}\ 4 \qquad\mathrm{(E)}\ 8$


	\item (\textit{AMC 12 12B 2004}) For some real numbers $a$ and $b$, the \href{https://artofproblemsolving.com/wiki/index.php?title=Equation}{equation}

 \[8x^3 + 4ax^2 + 2bx + a = 0\]
has three distinct positive roots. If the sum of the base-$2$ \href{https://artofproblemsolving.com/wiki/index.php?title=Logarithm}{logarithms} of the roots is $5$, what is the value of $a$? 


 $\mathrm{(A)}\ -256 \qquad\mathrm{(B)}\ -64  \qquad\mathrm{(C)}\ -8 \qquad\mathrm{(D)}\ 64  \qquad\mathrm{(E)}\ 256$


	\item (\textit{AMC 12 12B 2004}) Points $A$ and $B$ are on the parabola $y=4x^2+7x-1$, and the origin is the midpoint of $AB$. What is the length of $AB$?


 $\mathrm{(A)}\ 2\sqrt5 \qquad \mathrm{(B)}\ 5+\frac{\sqrt2}{2} \qquad \mathrm{(C)}\ 5+\sqrt2 \qquad \mathrm{(D)}\ 7 \qquad \mathrm{(E)}\ 5\sqrt2$


	\item (\textit{AMC 12 12B 2004}) A truncated \href{https://artofproblemsolving.com/wiki/index.php?title=Cone}{cone} has horizontal bases with radii $18$ and $2$. A \href{https://artofproblemsolving.com/wiki/index.php?title=Sphere}{sphere} is tangent to the top, bottom, and lateral surface of the truncated cone. What is the \href{https://artofproblemsolving.com/wiki/index.php?title=Radius}{radius} of the sphere? 


 $\mathrm{(A)}\ 6 \qquad\mathrm{(B)}\ 4\sqrt{5} \qquad\mathrm{(C)}\ 9 \qquad\mathrm{(D)}\ 10 \qquad\mathrm{(E)}\ 6\sqrt{3}$


	\item (\textit{AMC 12 12B 2004}) Each face of a \href{https://artofproblemsolving.com/wiki/index.php?title=Cube}{cube} is painted either red or blue, each with probability $1/2$. The color of each face is determined independently. What is the probability that the painted cube can be placed on a horizontal surface so that the four vertical faces are all the same color?


 $\textbf{(A)}\ \frac14 \qquad \textbf{(B)}\ \frac {5}{16} \qquad \textbf{(C)}\ \frac38 \qquad \textbf{(D)}\ \frac {7}{16} \qquad \textbf{(E)}\ \frac12$


	\item (\textit{AMC 12 12B 2004}) The graph of $2x^2 + xy + 3y^2 - 11x - 20y + 40 = 0$ is an ellipse in the first quadrant of the $xy$-plane. Let $a$ and $b$ be the maximum and minimum values of $\frac yx$ over all points $(x,y)$ on the ellipse. What is the value of $a+b$?


 $\mathrm{(A)}\ 3 \qquad\mathrm{(B)}\ \sqrt{10} \qquad\mathrm{(C)}\ \frac 72 \qquad\mathrm{(D)}\ \frac 92 \qquad\mathrm{(E)}\ 2\sqrt{14}$


	\item (\textit{AMC 12 12B 2004}) The square


 $\begin{tabular}{|c|c|c|} \hline 50 & \textit{b} & \textit{c} \\ \hline \textit{d} & \textit{e} & \textit{f} \\ \hline \textit{g} & \textit{h} & 2 \\ \hline \end{tabular}$

 is a multiplicative magic square. That is, the product of the numbers in each row, column, and diagonal is the same. If all the entries are positive integers, what is the sum of the possible values of $g$?


 $\textbf{(A)}\ 10 \qquad \textbf{(B)}\ 25 \qquad \textbf{(C)}\ 35 \qquad \textbf{(D)}\ 62 \qquad \textbf{(E)}\ 136$


	\item (\textit{AMC 12 12B 2004}) The \href{https://artofproblemsolving.com/wiki/index.php?title=Polynomial}{polynomial} $x^3 - 2004 x^2 + mx + n$ has \href{https://artofproblemsolving.com/wiki/index.php?title=Integer}{integer} coefficients and three distinct positive zeros. Exactly one of these is an integer, and it is the sum of the other two.  How many values of $n$ are possible?


 $\mathrm{(A)}\ 250,000 \qquad\mathrm{(B)}\ 250,250 \qquad\mathrm{(C)}\ 250,500 \qquad\mathrm{(D)}\ 250,750 \qquad\mathrm{(E)}\ 251,000$


	\item (\textit{AMC 12 12B 2004}) In $\triangle ABC$, $AB = BC$, and $\overline{BD}$ is an \href{https://artofproblemsolving.com/wiki/index.php?title=Altitude}{altitude}. Point $E$ is on the extension of $\overline{AC}$ such that $BE = 10$. The values of $\tan \angle CBE$, $\tan \angle DBE$, and $\tan \angle ABE$ form a \href{https://artofproblemsolving.com/wiki/index.php?title=Geometric\_progression}{geometric progression}, and the values of $\cot \angle DBE,$ $\cot \angle CBE,$ $\cot \angle DBC$ form an \href{https://artofproblemsolving.com/wiki/index.php?title=Arithmetic\_progression}{arithmetic progression}. What is the area of $\triangle ABC$?


 
                    \begin{figure}[H]
                        \centering
                        \includegraphics[width=0.5\textwidth]{Problem/AMC_12_12B_2004_24_Problem_Figure_1.png}
                    \end{figure}
                    
$\mathrm{(A)}\ 16 \qquad\mathrm{(B)}\ \frac {50}3 \qquad\mathrm{(C)}\ 10\sqrt{3} \qquad\mathrm{(D)}\ 8\sqrt{5} \qquad\mathrm{(E)}\ 18$


	\item (\textit{AMC 12 12B 2004}) Given that $2^{2004}$ is a $604$-\href{https://artofproblemsolving.com/wiki/index.php?title=Digit}{digit} number whose first digit is $1$, how many \href{https://artofproblemsolving.com/wiki/index.php?title=Element}{elements} of the \href{https://artofproblemsolving.com/wiki/index.php?title=Set}{set} $S = \{2^0,2^1,2^2,\ldots ,2^{2003}\}$ have a first digit of $4$? 


 $\mathrm{(A)}\ 194 \qquad\mathrm{(B)}\ 195 \qquad\mathrm{(C)}\ 196 \qquad\mathrm{(D)}\ 197 \qquad\mathrm{(E)}\ 198$


\end{enumerate}